\section{Space Complexity --- Constructible}
%%% generated by the blueprint pipeline from TCSlib.Complexity.SpaceComplexity.Constructible
\subsection*{Space Complexity --- Constructible: Introduction}

Declarations of the Lean module \texttt{TCSlib.Complexity.SpaceComplexity.Constructible}.

\subsection*{Space Complexity --- Constructible: Results}

\begin{definition}[Space-constructible functions]
\lean{Complexity.SpaceConstructible}
\statementsource{Sav70}{pdf-c8033d2fd705-p003-b003}
\label{Complexity.SpaceConstructible}
\leanok
\uses{Complexity.logSpace, Turing.FinTM, Turing.FinTM.ComputesInSpace}
A function $S:\mathbb{N}\to\mathbb{N}$ is space constructible if $\lfloor\log_2
n\rfloor+1\le S(n)$ for every natural number $n$, and there exist a positive natural
number $c$ and a finite Boolean Turing machine that, on every Boolean input $x$, outputs
the binary representation of $S(|x|)$ using at most $cS(|x|)$ work-tape cells.
\end{definition}

\begin{lemma}[Length of the binary expansion]
\lean{Complexity.b3a_bits_length}
\label{Complexity.b3a_bits_length}
\leanok
\uses{Complexity.LogProg.length_bits_le_log, Complexity.logSpace}
\difficulty{3}
For every nonzero natural number $n$, the length of its binary digit list equals
$\log_2(n)+1$, where $\log_2(n)$ denotes the integer logarithm used in the definition of
logarithmic space.
\end{lemma}

\begin{definition}[B3a Word]
\lean{Complexity.b3a_word}
\label{Complexity.b3a_word}
\leanok
For a natural number $n$, the associated Boolean word is the list whose first entry is
true and whose remaining entries are the binary digits of $n$ after its first digit is
removed.
\end{definition}

\begin{lemma}[Length of the B3a word]
\lean{Complexity.b3a_word_length}
\label{Complexity.b3a_word_length}
\leanok
\uses{Complexity.b3a_bits_length, Complexity.b3a_word, Complexity.logSpace}
\difficulty{4}
For every natural number $n$, the word formed by placing a true bit before the tail of
the binary expansion of $n$ has length $\lfloor\log_2 n\rfloor+1$.
\end{lemma}

\begin{lemma}[Restricting buffer tape contents]
\lean{Complexity.b3a_buffer_restrict}
\label{Complexity.b3a_buffer_restrict}
\leanok
\uses{Turing.FinTM.bufferTape}
\difficulty{4}
For every Boolean word $w$ and integer cell position $z$, restricting the buffer's
contents to positions satisfying $0\le z<|w|$ leaves its value at $z$ unchanged.
\end{lemma}

\begin{lemma}[Overwrite the first buffer cell]
\lean{Complexity.b3a_buffer_overwrite}
\label{Complexity.b3a_buffer_overwrite}
\leanok
\uses{Turing.FinTM.bufferTape}
\difficulty{4}
For every Boolean word $w$ and integer cell position $z$, the tape obtained from the
buffer for $w$ by writing true at cell $0$ has the same contents as the buffer storing
the word formed by replacing the first symbol of $w$ with true (or using the one-symbol
word true when $w$ is empty).
\end{lemma}

\begin{definition}[Initial configuration with stored word]
\lean{Complexity.b3a_cfg}
\label{Complexity.b3a_cfg}
\leanok
\uses{Turing.Cfg, Turing.Cfg.init, Turing.FinTM.bufferTape}
Let $k$ be a positive number of work tapes, let $S$ be a state set, and let $s\in S$ and
$x,w$ be Boolean words. Define a configuration on input $x$ in state $s$ with output
empty, with the first work tape initialized to contain $w$, and with every other work
tape blank.
\end{definition}

\begin{lemma}[State renaming of the initial buffer configuration]
\lean{Complexity.b3a_cfg_map}
\label{Complexity.b3a_cfg_map}
\leanok
\uses{Complexity.b3a_cfg, Turing.Cfg.mapState}
\difficulty{1}
Let $k$ be a positive integer, let $S,T$ be sets, let $f:S\to T$, let $s\in S$, and let
$x,w$ be Boolean words. Renaming the state of the initial $k$-tape configuration on
input $x$, with $w$ stored on the first work tape, by $f$ gives exactly the
corresponding configuration with initial state $f(s)$.
\end{lemma}

\begin{lemma}[Embedding the initial buffer configuration]
\lean{Complexity.b3a_embed_cfg}
\label{Complexity.b3a_embed_cfg}
\leanok
\uses{Complexity.b3a_cfg, Turing.MultiTapeTM, Turing.embedEmitCfg, Turing.embedEmitCfg_selected_pos, Turing.embedEmitCfg_selected_tape, Turing.embedEmitTM_frame}
\difficulty{5}
Let $m,k$ be positive integers, let $S$ be a set, let $R$ be a machine with $m$ work
tapes and Boolean symbols, and let $\iota$ be an injection from the $m$ tape indices
into the $k$ tape indices such that $\iota(0)=0$. For any $s\in S$ and Boolean words
$x,w$, embedding the configuration with initial state $s$ and $w$ on its first tape into
the $k$-tape system, with all unused tapes blank, head positions zero, and empty output,
gives the corresponding $k$-tape configuration.
\end{lemma}

\begin{lemma}[Space bound for an embedded machine]
\lean{Complexity.b3a_embed_space}
\label{Complexity.b3a_embed_space}
\leanok
\uses{Turing.Cfg, Turing.MultiTapeTM, Turing.MultiTapeTM.spaceUsed, Turing.MultiTapeTM.spaceUsedByTape_le_spaceUsed, Turing.embedEmitCfg, Turing.embedEmitTM, Turing.embedEmitTM_visitedByTapeHead, Turing.embedEmitTM_visitedByTapeHead_frame}
\difficulty{4}
Let $m,k$ be natural numbers, let $S$ be a set, let $R$ be a machine with $m$ work tapes
and Boolean symbols, let $\iota$ inject the $m$ tape indices into the $k$ tape indices,
let $d$ be a configuration of $R$ on a Boolean input $x$, and let $t\in\mathbb N$. The
machine obtained by embedding $R$, started from the configuration that embeds $d$ with
every unused tape blank and its head at the origin, uses at most
$k\bigl(\operatorname{space}_R(d,t)+1\bigr)$ total space through time $t$.
\end{lemma}

\begin{lemma}[Space bound for sequential machine composition]
\lean{Complexity.b3a_seam_space}
\label{Complexity.b3a_seam_space}
\leanok
\uses{Turing.Cfg, Turing.Cfg.mapState, Turing.MultiTapeTM, Turing.MultiTapeTM.runFrom, Turing.MultiTapeTM.spaceUsed, Turing.MultiTapeTM.spaceUsedByTape, Turing.MultiTapeTM.visitedByTapeHead, Turing.seamCompTM, Turing.seamCompTM_visitedByTapeHead_ofCfg}
\difficulty{3}
Let $S,T$ be sets, let $k\in\mathbb N$, and let $R,Q$ be machines with $k$ work tapes
and Boolean symbols, with exit state $e_x\in S$ for $R$ and entry state $e_y\in T$ for
$Q$. Suppose configurations $d,e$ of $R$ on a Boolean input satisfy $R^a(d)=e$, the
state of $e$ is $e_x$, and no earlier run $R^u(d)$ for $u<a$ has state $e_x$. Then the
machine that runs $R$ until its exit state and switches to $Q$ at entry state $e_y$,
started from $d$ with its states relabeled into the left side of $S\oplus T$, uses at
most $\operatorname{space}_R(d,a)+\operatorname{space}_Q(e',b)$ space through time
$a+1+b$, for every $b\in\mathbb N$, where $e'$ is $e$ with its state replaced by $e_y$.
\end{lemma}

\begin{theorem}[A machine for writing the binary length]
\lean{Complexity.b3a_counter}
\label{Complexity.b3a_counter}
\leanok
\uses{Complexity.b3a_buffer_restrict, Complexity.b3a_cfg, Turing.MultiTapeTM, Turing.MultiTapeTM.runFrom, Turing.MultiTapeTM.spaceUsed, Turing.MultiTapeTM.spaceUsedByTape, Turing.exists_inputLengthTM}
\difficulty{5}
There exist a natural number $K$, a finite state set $S$ with decidable equality, a
deterministic machine $R$ with $K+1$ work tapes and Boolean symbols, entry and exit
states in $S$, and a natural number $B$, such that for every Boolean word $x$ there is a
time $t$ when $R$, started with the empty word on its first work tape and $x$ as input,
is in the exit configuration with the binary expansion of $|x|$ on that tape. It has not
reached the exit state at any earlier time $u<t$, and its total space use through time
$t$ is at most $B$ times one plus the length of the binary expansion of $|x|$.
\end{theorem}

\begin{theorem}[Overwriting the first buffer symbol]
\lean{Complexity.b3a_overwrite}
\label{Complexity.b3a_overwrite}
\leanok
\uses{Complexity.b3a_buffer_overwrite, Complexity.b3a_cfg, Turing.Cfg.init, Turing.FinTM.bufferTape, Turing.MultiTapeTM.runFrom, Turing.MultiTapeTM.spaceUsed, Turing.writeConstTM, Turing.writeConstTM_run_ofCfg}
\difficulty{5}
For any Boolean words $x,w$, the specified one-tape machine, started in its writing
phase with $w$ on its work tape, reaches after exactly four steps the done configuration
with work-tape content $\mathrm{true}::\mathrm{tail}(w)$. Its state is not done at any
earlier time $u<4$, and its total space use through step four is at most $2$.
\end{theorem}

\begin{definition}[Virtual coordinate embedding]
\lean{Complexity.b3a_virtualSelect}
\label{Complexity.b3a_virtualSelect}
\leanok
For natural numbers $K,m$, define an injection from the $m+1$ coordinates indexed by
$0,1,\ldots,m$ into the $K+m+1$ coordinates indexed by $0,1,\ldots,K+m$: it sends $0$ to
$0$ and each $i\ge1$ to $K+i$.
\end{definition}

\begin{lemma}[Virtual configuration as a buffered configuration]
\lean{Complexity.b3a_virtual_cfg}
\label{Complexity.b3a_virtual_cfg}
\leanok
\uses{Complexity.b3a_cfg, Turing.Cfg.init, Turing.FinTM, Turing.FinTM.bufferTape, Turing.MultiTapeTM.initCfg, Turing.vhostCfg}
\difficulty{4}
Let $L$ be a finite-state Turing machine over Boolean symbols, and let $x,y$ be Boolean
words. The virtual configuration formed from the initial configuration of $L$ on input
$y$, with boundary tag true, virtual input-head position $1$ on input $x$, and empty
output prefix, is the configuration on input $x$ whose state is the initial state of $L$
paired with true, and whose first work tape stores $y$ while every other work tape is
blank.
\end{lemma}

\begin{theorem}[A logarithmic-space machine for binary lengths]
\lean{Complexity.b3a_log_machine}
\label{Complexity.b3a_log_machine}
\leanok
\uses{Complexity.LogProg.length_bits_le_log, Complexity.b3a_cfg, Complexity.b3a_cfg_map, Complexity.b3a_counter, Complexity.b3a_embed_cfg, Complexity.b3a_embed_space, Complexity.b3a_overwrite, Complexity.b3a_seam_space, Complexity.b3a_virtualSelect, Complexity.b3a_virtual_cfg, Complexity.b3a_word, Complexity.b3a_word_length, Complexity.logSpace, Turing.Cfg.init, Turing.Cfg.mapState, Turing.FinTM, Turing.FinTM.ComputesInSpace, Turing.FinTM.VirtualTag, Turing.FinTM.computesFunInTime_lengthBits_spaceUsed, Turing.FinTM.computesInTime_iff, Turing.MultiTapeTM.initCfg, Turing.MultiTapeTM.runFrom, Turing.MultiTapeTM.spaceUsed, Turing.WritePhase, Turing.embedEmitCfg, Turing.embedEmitTM, Turing.embedEmitTM_runFrom, Turing.length_bits_le_self, Turing.seamCompTM, Turing.seamCompTM_firstReturn_ofCfg, Turing.seamCompTM_run_ofCfg, Turing.vhostCfg, Turing.vhostEmitTM, Turing.vhostEmitTM_runFrom, Turing.vhostEmitTM_spaceUsed_le, Turing.writeConstTM}
\difficulty{7}
There exist a positive natural number $A$ and a finite-state machine $M$ over Boolean
symbols such that, for every Boolean word $x$, $M$ halts with output the binary
expansion of $\lfloor\log_2 |x|\rfloor+1$ and uses at most $A(\lfloor\log_2
|x|\rfloor+2)$ work-tape cells. Here the logarithmic expression is understood with the
convention that its value is $1$ when $|x|=0$.
\end{theorem}

\begin{theorem}[Space constructibility of logarithmic space]
\lean{Complexity.spaceConstructible_logSpace}
\statementsource{Sav70}{pdf-c8033d2fd705-p003-b003}
\label{Complexity.spaceConstructible_logSpace}
\leanok
\uses{Complexity.SpaceConstructible, Complexity.b3a_log_machine, Complexity.logSpace}
\difficulty{4}
Define $L(n)=1+\lfloor\log_2 n\rfloor$ for $n\ge 1$, and set $L(0)=1$. There are a
positive integer $c$ and a finite binary Turing machine that, on every binary input $x$,
outputs the binary representation of $L(|x|)$ and uses at most $cL(|x|)$ work-tape
cells; moreover, $L(n)\le L(n)$ for every $n\in\mathbb N$.
\end{theorem}

\begin{theorem}[Linear space constructibility]
\lean{Complexity.spaceConstructible_linear}
\label{Complexity.spaceConstructible_linear}
\leanok
\uses{Complexity.SpaceConstructible, Turing.FinTM.bufferedCompTM, Turing.FinTM.bufferedCompTM_computesInTime, Turing.FinTM.computesFunInTime_lengthBits_spaceUsed, Turing.FinTM.computesFunInTime_prepend_spaceUsed, Turing.MultiTapeTM.initCfg, Turing.MultiTapeTM.spaceUsed, Turing.MultiTapeTM.spaceUsed_le_mul_succ}
\difficulty{4}
The function $S(n)=n+1$ is space constructible: for every $n\in\mathbb{N}$,
$\lfloor\log_2 n\rfloor+1\le n+1$. Moreover, there are a positive integer $c$ and a
finite Turing machine over the Boolean alphabet that, on every Boolean input $x$,
outputs the binary digit list of $|x|+1$ and uses at most $c(|x|+1)$ work-tape cells.
\end{theorem}
